\section{System Design and the Provenance Gap}
\label{sec:system}

The preceding sections evaluate models. A deployed system is a different object,
and the distance between the two is where most claims about this work must stop.
We describe what the accompanying deployment project (AeroShield) contains, and
then report a finding that belongs to this paper's central argument even though it
concerns engineering rather than detection.

\paragraph{Architecture, as design and not as measurement.} The project targets a
Jetson Nano ($4$\,GB, Maxwell generation, JetPack 4.6 / L4T 32.7.6, CUDA 10.2,
TensorRT 8.2, Python 3.6). Inference is onboard and the ground station is a
client. A FastAPI service backed by a PostGIS store handles missions, detections,
operator adjudication and report generation, with a web frontend providing a map
view, a Grad-CAM comparison and an adjudication workflow.

Two features of this design matter for how the present results should be read.

First, \emph{the backend runs no model}. Its image-upload detection endpoint is a
development stub returning a fixed detection with a fixed confidence, and the
endpoint's own docstring cites the project requirements (\S2.2) as ruling out
server-side inference. This is the right architecture for the task---the sensor
is on the aircraft, and the ground station's job is bookkeeping, human review and
reporting---but it has a consequence for reporting: no throughput, latency or
accuracy figure obtained from the backend describes detection performance. A
system diagram that shows a detector on both sides of the radio link is
architecturally wrong as well as unsupported.

Second, \emph{the export path is constrained by the device}. The detector is
exported to ONNX at a fixed $640\times640$ input with opset 12 and non-maximum
suppression left out of the graph, then compiled on the target with
\texttt{trtexec --fp16}. The opset ceiling is set by the platform, not by
preference: JetPack 4.6 ships TensorRT 8.2, and an engine produced for a newer
TensorRT cannot be loaded. The engine must be built on the device; an engine
compiled on a desktop GPU will not run on Maxwell. The ONNX file is portable and
the engine is not, which is a common source of edge numbers that cannot be
reproduced from the repository that reports them.

\paragraph{The provenance gap.} Our evaluation and the deployment diverged, and
neither repository records where. The models evaluated here are YOLOv8s
($11.2$\,M parameters) and YOLOv8n ($3.0$\,M) trained on the rebuilt split
$3{,}907/181/189$ and evaluated on a leak-audited test set. The deployment
project's committed evidence points elsewhere: its training script sets the base
model to \texttt{yolov8m.pt}, its committed training report documents a
$75$-epoch YOLOv8m run on the pre-rebuild split ($5{,}723/455/221$), and the ONNX
file it ships is $103.6$\,MB beside a $42.7$\,MB PyTorch checkpoint of the same
name. The parameter ratio between the two backbones is $25.9/11.2 = 2.31$; the
file-size ratio is $103.6/42.7 = 2.32$. Both artifacts declare one class
(\texttt{Landmine}), so this is not the class-count confusion we corrected
earlier in this project's own documentation---it is a question of which backbone
and which split. Which file was actually flashed to the device is not recorded in
either repository.

\begin{table}[t]
\centering
\caption{What was evaluated against what was deployed. The right-hand column is
assembled from the deployment project's training script, its committed report and
the size of the ONNX file it ships; the three agree with each other. Neither
repository records which file reached the device.}
\label{tab:artifact}
\scriptsize
\begin{tabular}{lll}
\toprule
 & evaluated here & deployment artifacts \\
\midrule
backbone & YOLOv8s & YOLOv8m \\
parameters & $11.2$\,M & $25.9$\,M \\
classes & 1 (\texttt{Landmine}) & 1 (\texttt{Landmine}) \\
split & v2, $3{,}907/181/189$ & v1, $5{,}723/455/221$ \\
leak-audited & yes, $30$ removed & no \\
provenance & manifest digest & not recorded \\
evaluated here & --- & no \\
\bottomrule
\end{tabular}
\end{table}

This is a leakage failure of a different kind, and it is worth naming precisely,
because the vocabulary of the field does not cover it. The literature on leakage
describes \emph{data} crossing a split boundary. What happened here is an
\emph{artifact} crossing a project boundary: a checkpoint was produced, reported,
and deployed, and the report describing it was filed in a directory that was later
emptied by a rebuild. Every individual step was reasonable. The training run is
documented; the report is committed; the weights directory has a README
explaining how to regenerate the checkpoint. The failure is compositional. The
only link between the evaluated model and the deployed model was a filename, and
the filename was reused.

The methodological recommendation is narrow and cheap. A commit hash identifies
the \emph{code} that produced a model, not the model. A checkpoint manifest---file
name, SHA-256, backbone, class list, split digest, and the export parameters used
to produce the ONNX---identifies the model, and it can be verified from the
artifact in a few seconds without re-running anything. This paper ships such a
manifest for the evaluated models (Section~\ref{sec:availability}). The deployment
side does not, and the cost of that absence is exactly the ambiguity in
Table~\ref{tab:artifact}: a defensible claim that a detector of this family runs
on this hardware, and no defensible claim about which one.

\paragraph{What remains unmeasured.} The deployment repository contains no
committed timing record. Its edge numbers appear in prose in its README, no run
artifacts are stored, and therefore none of them can be verified, re-derived, or
compared with Table~\ref{tab:size}. We do not import them. The gap is not a
contradiction to be argued with but an absence to be filled: the models evaluated
here are exported and ready to run onboard, and producing an on-device number for
one of them is an instrumentation task---build the engine on the device, run the
existing per-stage timing harness, commit its summary---rather than a change of
method. Until that is done this paper makes no claim about edge throughput,
energy, or thermal behaviour, and the deployment project's architecture should be
read as design intent rather than as a validated result.